\section{Introduction}
\label{sec:introduction}

Gameplay glitches are visible software failures that disrupt interaction, presentation,
or progression. They are usually temporal events: a character sticks in place, an animation
freezes, a projectile behaves inconsistently, or motion violates the expected game dynamics.
This makes gameplay glitch detection different from static screenshot anomaly detection.
Benchmarks such as GlitchBench, TempGlitch, and VideoGlitchBench make this distinction visible:
some tasks ask whether a single frame looks unusual, while temporal glitch detection asks whether
an episode evolves in an unexpected way~\cite{glitchbench,tempglitch,videoglitchbench}.

The core insight of this paper is that a gameplay glitch is a violation of expected game
dynamics. When a character teleports unexpectedly, a projectile freezes, or an avatar becomes
stuck, the observed transition no longer matches the latent dynamics that normal gameplay would
have implied. This maps naturally onto world-model surprise: if a model has learned only normal
gameplay, a glitch should manifest as an anomalously large prediction error in latent space.
Unlike pixel-level anomaly detection, latent-space surprise is less sensitive to superficial
visual variation such as camera angle, lighting, or texture detail that may change during normal
play but do not themselves indicate a bug.

Joint-Embedding Predictive Architectures (JEPAs) are a particularly natural fit for this
setting because they predict in representation space rather than reconstructing pixels directly,
focusing the objective on semantically meaningful dynamics instead of every visual detail
in the frame~\cite{lecun2022path,ijepa,vjepa,vjepa2}. LeWorldModel extends this family to
pixel-based world modeling and gives us an artifact-backed gameplay backbone for testing whether
latent prediction error can serve as a label-free glitch signal~\cite{lewm}. In this paper we
use that JEPA-style backbone as an unsupervised detector: train on normal gameplay, score each
window by latent prediction error, aggregate to the episode level, and calibrate thresholds using
normal-only episodes.

This paper asks a bounded empirical question rather than introducing a new large model:
does JEPA-style latent surprise produce useful glitch-detection evidence under leakage-aware
protocols on public gameplay benchmarks? We answer that question with one positive same-support
TempGlitch comparison, one stronger secondary frozen R5-XGame result family, and two honest
negative analyses that clarify where temporal latent surprise does and does not help.

The contributions of this paper are:
\begin{itemize}[leftmargin=*]
\item \textbf{Method}: LeWM-Glitch, a JEPA-based latent prediction error (LPE) signal
for label-free gameplay glitch detection, instantiated from artifact-backed LeWorldModel
checkpoints with source-, pair-, and episode-aware role separation before windowing~\cite{lewm};
\item \textbf{Multi-benchmark leakage-aware evaluation}: a frozen pair-disjoint TempGlitch
follow-up and a complementary R5-XGame evaluation, both under strict normal-only calibration
protocols, establishing LeWM-Glitch's performance on two independently bounded
splits~\cite{tempglitch,videoglitchbench};
\item \textbf{Honest negative findings}: a bounded GlitchBench K2 negative result (AUROC 0.5 for
the current LeWM configuration vs.\ AUROC 1.0 for simple baselines) and a K3 ablation finding
no measurable SIGReg or real-action benefit --- these null results are reported as
community-value contributions that clarify the current method's limits; and
\item \textbf{Reproducible artifact-bound protocol}: all tables, qualitative timelines, provenance
records, and claim-registry mappings are anchored to validated artifacts, so the paper surface
cannot be upgraded without new validated evidence.
\end{itemize}

The scope is deliberately careful. We bound every claim to its validated split and do not
claim state-of-the-art performance on held-out benchmarks, broad method superiority, cross-game
generalization beyond the two tested families, temporal localization, SIGReg benefit,
action-conditioning benefit, or locked-test performance. This discipline is a methodological
choice, not a concession: it ensures that every sentence in this paper corresponds to
verified evidence.
